% 巨行星（综述）
% license CCBYSA3
% type Wiki

本文根据 CC-BY-SA 协议转载翻译自维基百科\href{https://en.wikipedia.org/wiki/Giant_planet}{相关文章}。

巨行星（giant planet），有时也被称为类木行星（jovian planet）（Jove 是罗马神话中木星〔Jupiter〕的另一名称），是一类比地球大得多、类型多样的行星。巨行星通常主要由低沸点物质（挥发物）构成，而非岩石或其他固态物质，但“超级地球”（mega-Earths）这一类别也确实存在。在太阳系中共有四颗此类行星：木星、土星、天王星和海王星。此外，已有大量系外巨行星被发现。

巨行星有时被称为气体巨行星（gas giants），但许多天文学家现在将此称谓仅用于木星和土星，并将成分不同的天王星与海王星归类为冰巨行星（ice giants）。然而，这两种名称都可能导致误解：在太阳系中，所有巨行星主要由处于临界点之上的流体组成，在此状态下气相与液相不再具有明确区分。木星和土星主要由氢与氦构成，而天王星与海王星则由水、氨和甲烷构成。

关于低质量褐矮星与高质量气体巨行星（约 \(13\,M_{\mathrm{J}}\)）之间的定义差异仍存在争议。一种观点基于行星的形成过程；另一种观点则基于行星内部物理性质。争论的一部分涉及褐矮星是否必须在其演化史中的某一阶段经历过核聚变，这一点是否应被视为定义它们的必要条件。\(^\text{[1]}\)
